% ==========================================================================
% INFO 4617 Web Data Science -- Week 8: Revise a pull request
% Build: `make` in handouts/, or `latexmk -pdf revise-a-pull-request.tex`
% here. Every screen is GitHub's, seen signed in, so every figure is a hand
% capture of the demonstration pull request, textbook #228 (img/IMAGES.md):
% \handcapture prints it once img/NAME.png exists, and leaves it out until
% then.
% ==========================================================================
\makeatletter\def\input@path{{../common/}{handouts/common/}}\makeatother
\documentclass{handout}

\title[INFO 4617 · Week 8 · Revise a pull request]{Revise a pull request}
\subtitle{Week 8 · a reference, and a Friday activity · about 20 minutes}
\author{Brian C. Keegan, Ph.D.}
\institute{Department of Information Science, University of Colorado Boulder}
\date{October 9, 2026}

% A numbered figure: [width]{path from the repo root}{alt text}{caption}
\newcounter{handoutfig}
\newcommand{\activityfigure}[4][\linewidth]{%
  \par\smallskip\refstepcounter{handoutfig}%
  {\centering
   \BeginAccSupp{method=pdfstringdef,Alt={#3}}%
   \includegraphics[width=#1]{#2}%
   \EndAccSupp{}\par}%
  \smallskip{\small\raggedright\textbf{Figure \thehandoutfig.} #4\par}}
% A hand capture: {file in handouts/week-08/img/}{alt text}{caption}
\newcommand{\handcapture}[3]{%
  \IfFileExists{handouts/week-08/img/#1}%
    {\activityfigure{handouts/week-08/img/#1}{#2}{#3}}%
    {\IfFileExists{../../handouts/week-08/img/#1}%
      {\activityfigure{handouts/week-08/img/#1}{#2}{#3}}{}}}

% Text a conflict editor shows: the code block's look, without Python's colors
\lstnewenvironment{plain}{\lstset{style=handoutcode,language={}}}{}

\setlist[enumerate]{leftmargin=*,itemsep=1.5pt,topsep=2pt}
\setlist[itemize]{leftmargin=*,itemsep=1pt,topsep=2pt}

\begin{document}
\maketitle

\begin{frame}{What you do}
  \begin{columns}[T]
    \begin{column}{0.57\textwidth}
      A review of your pull request asks for changes. You make them on the
      same pull request. A pull request follows its branch: each commit you
      add to the branch shows in the pull request at once, next to the
      comments it answers. A new pull request would lose that history, and
      the reviewer would have to start again.

      You do all of it in the browser, with one of three tools. The table
      on the right says which tool fits which kind of comment.

      The example on each page is textbook pull request \#228, a
      demonstration that the instructor set up. Its base is a copy of
      \texttt{main} named \texttt{demo/week08-base}. On your pull request,
      the base is \texttt{main}.
    \end{column}
    \begin{column}{0.39\textwidth}
      \begin{block}{Which tool for which comment}
        \small
        \begin{tabular}{@{}>{\raggedright\arraybackslash}p{0.46\linewidth}%
                        >{\raggedright\arraybackslash}p{0.46\linewidth}@{}}
          A comment with a \textbf{suggestion} box & Step 2: \textbf{Commit suggestion}\\ \addlinespace[4pt]
          A comment that asks for a change in words & Step 3: \textbf{Edit file}\\ \addlinespace[4pt]
          The notice \textit{This branch has conflicts} & Step 4: \textbf{Resolve conflicts}
        \end{tabular}
      \end{block}
      \begin{alertblock}{Do not}
        \small
        \begin{itemize}
          \item open a new pull request, or close this one;
          \item click \textbf{Merge pull request};
          \item change a file in \texttt{notebooks/} or \texttt{images/};
          \item make changes that the review did not ask for.
        \end{itemize}
      \end{alertblock}
      \begin{exampleblock}{In class}
        \small Work in pairs, on one partner's pull request. One partner
        works in the browser, and the other reads this handout aloud. Then
        change roles.
      \end{exampleblock}
    \end{column}
  \end{columns}
\end{frame}

\begin{frame}{1 · Read the whole review}
  \begin{enumerate}
    \item Sign in to GitHub, and open your pull request
      (\url{https://github.com/cuinfoscience/Web-Data-Science-Book/pulls}).
    \item On the \textbf{Conversation} tab, read the review's summary, then
      each comment under it. A comment on a line shows the line above it.
    \item Make a list: one item for each change the review asks for, and
      the tool from the table on page~1. \#228's review asks for two: fix a
      typo (a suggestion), and add \texttt{-{}-verbose} to the
      \texttt{curl} line (a comment). Its conflict is a third item.
    \item If you don't understand a comment, or you disagree, reply in its
      thread before you change anything. Say what you think, and why.
  \end{enumerate}
\end{frame}

\begin{frame}{2 · A suggestion: Commit suggestion}
  A reviewer who writes the exact new words puts them in a
  \textbf{suggestion} box. GitHub shows it as the old line in red and the
  new line in green, with a button under it.
  \begin{enumerate}
    \item Check the green line. It replaces the red line exactly as shown.
    \item Click \textbf{Commit suggestion}. Write a commit message that says
      what changed, for example \texttt{Ch. 5: fix the spelling of
      Network}. Click \textbf{Commit changes}.
    \item The comment now says the suggestion was committed, and the
      pull request shows one more commit.
  \end{enumerate}
  If a review has several suggestions, click \textbf{Add suggestion to
  batch} on each one, then \textbf{Commit suggestions} once. That makes one
  commit for all of them.
  \handcapture{pr-suggestion.png}{Screenshot of a review comment on pull
    request 228 with a suggestion box: the old line with the typo Netwrok in
    red, the new line with Network in green, and the Commit suggestion and
    Add suggestion to batch buttons under it.}{A suggestion on \#228: the
    old line in red, the new line in green, and \textbf{Commit suggestion}
    under them.}
\end{frame}

\begin{frame}{3 · A comment: Edit file on your branch}
  \begin{enumerate}
    \item Click \textbf{Files changed}. Find the file, and click its
      \textbf{$\cdots$} menu at the right of the file's name. Click
      \textbf{Edit file}.
    \item Check the top of the editor. It names your pull request's branch,
      not \texttt{main}. (\#228's branch is
      \texttt{demo/week08-revise}.) If it says \texttt{main}, or asks you
      to fork the repository, you opened the file from the wrong page: go
      back to your pull request's \textbf{Files changed}.
    \item Press \texttt{Ctrl+F} (\texttt{Cmd+F} on a Mac) in the editor, and
      find the line. Make the change that the comment asks for, and nothing
      else. Click \textbf{Preview}, and compare the red and green lines.
    \item Click \textbf{Commit changes…}. Write a commit message, and keep
      \textbf{Commit directly to the} \textit{your-branch} \textbf{branch}
      selected. Do not create a new branch: a new branch would need a new
      pull request. Click \textbf{Commit changes}.
  \end{enumerate}
  \handcapture{pr-edit-commit.png}{Screenshot of the Commit changes dialog
    after editing a file of pull request 228: a commit message, and the
    option Commit directly to the demo/week08-revise branch selected.}{The
    \textbf{Commit changes} dialog on \#228: the option that commits to the
    pull request's own branch is selected.}
\end{frame}

\begin{frame}{4 · A conflict: Resolve conflicts}
  The bottom of the \textbf{Conversation} tab says \textit{This branch has
  conflicts that must be resolved} when \texttt{main} changed the same lines
  as your pull request after you opened it. GitHub can't choose between the
  two versions, so you do. Resolve it when a review asks you to.
  \begin{enumerate}
    \item Click \textbf{Resolve conflicts}. GitHub opens an editor with
      each conflict between three marker lines:
  \end{enumerate}
\begin{plain}
<<<<<<< demo/week08-revise
- `curl` manual, for the `-v` flag this chapter uses: <https://curl.se/docs/manpage.html>
=======
- `curl` manual, for `-v` and `-I`: <https://curl.se/docs/manpage.html>
>>>>>>> demo/week08-base
\end{plain}
  \begin{enumerate}[resume]
    \item The first part is your branch's version; the second is the base's
      (on your pull request, \texttt{main}'s). Write the one line that keeps
      what each side meant. For \#228, the base added \texttt{-I}, and the
      review asked for \texttt{-{}-verbose}:
\begin{plain}
- `curl` manual, for `-v` (`--verbose`) and `-I`: <https://curl.se/docs/manpage.html>
\end{plain}
    \item Delete the three marker lines (\texttt{<{}<{}<}, \texttt{===},
      \texttt{>{}>{}>}). Then click \textbf{Mark as resolved}, and
      \textbf{Commit merge}.
  \end{enumerate}
  If the button is gray and says the conflict is \textit{too complex to
  resolve in the web editor}, stop. Comment on your pull request that you
  need help with a conflict, and the instructor will resolve it.
  \handcapture{pr-conflict-editor.png}{Screenshot of GitHub's conflict editor
    for pull request 228, with the three marker lines around the two versions
    of the curl line, and the Mark as resolved button.}{GitHub's conflict
    editor on \#228: the two versions of the line, between the marker
    lines, and \textbf{Mark as resolved}.}
\end{frame}

\begin{frame}{5 · Answer the review, and ask again}
  \begin{enumerate}
    \item Reply in each thread with what you did, for example
      \texttt{Done: added the long form.} Click \textbf{Resolve
      conversation} under a thread when its change is made. Leave a thread
      open if you disagreed and are waiting for an answer.
    \item In the right-hand sidebar, under \textbf{Reviewers}, click the
      circular arrows beside the reviewer's name, \textbf{Re-request
      review}. The reviewer gets a notice that your pull request is ready
      again.
    \item Watch the checks at the bottom of the \textbf{Conversation} tab.
      They run again after each commit. A red \textbf{Notebook sync} is
      normal for a change made in the browser. A red \textbf{Render} is
      yours to fix: click \textbf{Details}, find the first \texttt{ERROR}
      line, and fix it with step 3.
  \end{enumerate}
  \handcapture{pr-rerequest.png}{Screenshot of the Reviewers section of the
    pull request sidebar, with the pointer on the circular arrows button
    labeled Re-request review.}{\textbf{Re-request review}, the circular
    arrows beside the reviewer's name in the sidebar.}

  \begin{alertblock}{If your pull request came from your fork's main branch}
    \small
    The pull request shows \textit{wants to merge into
    cuinfoscience:main from your-name:main}. Then every commit you make on
    your fork's \texttt{main} joins this pull request. Commit only the
    review's changes there until it merges. Start your next change from a
    new branch. If your fork offers \textbf{Sync fork}, click
    \textbf{Update branch}, never \textbf{Discard commits}: discarding
    deletes your pull request's commits.
  \end{alertblock}
\end{frame}

\begin{frame}{When something goes wrong}\relax
  {\small
  \begin{tabular}{@{}>{\raggedright\arraybackslash}p{0.34\linewidth}%
                  >{\raggedright\arraybackslash}p{0.62\linewidth}@{}}
    \toprule
    \textbf{You see} & \textbf{Do this}\\
    \midrule
    \textbf{Commit suggestion} is gray, or the comment says
    \textit{Outdated} &
      The line changed after the comment. Make the change with
      \textbf{Edit file} (step 3).\\[3pt]
    The editor names \texttt{main}, or GitHub asks to fork the repository &
      You opened the file from the book, not from your pull request. Close
      the editor without committing, and start step 3 from \textbf{Files
      changed}.\\[3pt]
    Your commit opened a new pull request &
      You chose \textbf{Create a new branch}. Close the new pull request,
      with a comment that names your first one. Make the change again on the
      first pull request's branch.\\[3pt]
    \textbf{Resolve conflicts} is gray &
      The conflict is too complex for the web editor. Comment on your pull
      request, and the instructor will resolve it.\\[3pt]
    The checks say a workflow is \textit{awaiting approval} &
      A pull request from a fork waits until the instructor approves its
      checks. You need not do anything.\\[3pt]
    You closed your pull request by mistake &
      Click \textbf{Reopen pull request}, at the bottom of the
      \textbf{Conversation} tab. Its branch and commits are still there.\\
    \bottomrule
  \end{tabular}\par}

  \medskip
  \textbf{Next.} When the reviewer approves, the instructor merges your pull
  request after the Friday code review. Answer any new comments the same way.
\end{frame}
\end{document}
